\documentclass[letter,11pt]{article}
\usepackage{hyperref}
\usepackage{graphicx}
\usepackage{fancyhdr}
\usepackage{multicol}
\pagestyle{fancy}
\usepackage[letterpaper, margin=1in]{geometry}
\geometry{letterpaper}
\usepackage{parskip} % Disable initial indent

\usepackage[utf8]{inputenc}
\fancyhf{}
\renewcommand{\headrulewidth}{0pt} % Remove default underline from header package
\rhead{CYBR/CMSC 691 - Lab 10}
%\rhead{}
\lhead{\begin{picture}(0,0) \put(0,-10){\includegraphics[width=1.1cm]{Images/UMBC-vertical}} \end{picture}}
\cfoot{\thepage}
\rfoot{Spring 2020}
\lfoot{Lab 10}
%\AtEndDocument{\vfill \hfill \LaTeX}

\begin{document}

\section*{Lab 10: Malware $N$-Grams}
\paragraph{}Provided in the VM are a variety of malware \& goodware samples. Use the provided script to extract $n$-grams, and make a dataset based on those $n$-grams. The size of $n$-grams, number of top $n$-grams to keep, are user-selectable.

\textbf{Points:} 20

%\subsection*{Hints}
%\paragraph{}If you'd like, you may review code written in \href{https://www.golang.org}{Go} which does the same thing: \url{https://github.com/rjzak/gogrammer}. Creating an array of ngrams will take a lot of memory. use the hash of the ngram as the index into an array, and using the value at the array to increment a counter. Continue then using the $K$ most common $n$-grams.

\section*{Steps}
\begin{enumerate}
    \item Create a dataset based on the top 100 4-grams
    \item Create a dataset based on the top 500 4-grams
    \item Create a dataset based on the top 100 6-grams
    \item Create a dataset based on the top 500 6-grams
    \item Train \texttt{xgboost} models for each dataset.
    \item Use \texttt{shap} to see which features were more and less important.
\end{enumerate}

\section*{Questions}
\paragraph{}In a separate document, provide your answers to the following questions. The expected format is a Word document, or similar, with a few images demonstrating your progress. The assignment will be submitted via Blackboard.

\begin{itemize}
    \item How does the accuracy change between top 100 and 500 $n$-grams?
    \item How does the accuracy change between the use of 4 and 6-grams?
    \item Which features were more or less important?
    \item Speculate on what some of the features might \textit{mean}.
\end{itemize}

\end{document}